% titel: Blatt 1, A10: Gemischtypen / A11: Trennmethoden
% Teil II: der durchgerechnete Weg.

\subsection*{A10 --- Gemischtyp je Stoff}
\begin{itemize}
  \item \textbf{Sirup:} Lösung (Zucker in Wasser, homogen).
  \item \textbf{Beton (flüssig):} Suspension (Zementkörner in Wasser, heterogen).
  \item \textbf{Arzneitablette:} Gemenge (festes Pulver mehrerer Stoffe).
  \item \textbf{Mayonnaise:} Emulsion (Öl in Wasser, flüssig/flüssig).
  \item \textbf{Staub:} Rauch-ähnlich, feste Partikel in Gas (Aerosol).
  \item \textbf{Abgase:} meist Lösung/Gasgemisch (gasförmig homogen), teils
        mit Rauchanteil.
  \item \textbf{Schwarzpulver:} Gemenge (feste Bestandteile, heterogen).
\end{itemize}

\begin{ergebnis}
Sirup=Lösung, Beton=Suspension, Tablette=Gemenge, Mayonnaise=Emulsion,
Staub=Aerosol(fest/gas), Schwarzpulver=Gemenge.
\end{ergebnis}

\subsection*{A11 --- Geeignete Trennmethoden}
\begin{itemize}
  \item \textbf{Festes AgCl aus Lösung:} Filtration (Niederschlag bleibt am
        Filter zurück).
  \item \textbf{Farbstoffe aus zermahlener Blüte (in Aceton löslich):}
        Extraktion (mit Aceton herauslösen), ggf.\ Chromatographie zur
        Auftrennung der einzelnen Farbstoffe.
  \item \textbf{Methanol aus wässriger Lösung:} Destillation
        (unterschiedlicher Siedepunkt).
  \item \textbf{Schwefel (Smp.\ 114\,°C) neben Sand:} Extraktion (Schwefel
        in geeignetem Lösungsmittel) oder Destillation/Sublimation ---
        Schwefel schmilzt, Sand nicht, Trennüberführung über die
        Schmelztemperatur.
  \item \textbf{Gelöster Farbstoff (wasser-, in Benzin löslich):}
        Extraktion (in Benzin überführen, Schütteltrichter).
  \item \textbf{Blutplasma und Zellen:} Zentrifugation (Dichteunterschied).
\end{itemize}

\begin{kontrolle}
Meist mehrere Wege sinnvoll; die Auflösung in eine passende Phase
(Extraktion) bzw.\ Phasentrennung (Filtration/Zentrifugation/Destillation) ist
das Leitprinzip.
\end{kontrolle}

\begin{merksatz}
Trennregel: \emph{Phase und Löslichkeit} bestimmen die Methode ---
unlöslich fest/flüssig $\rightarrow$ filtrieren/zentrifugieren; flüssig/flüssig mit
unterschiedl. Sdp. $\rightarrow$ destillieren; gelöst übertragbar $\rightarrow$ extrahieren.
\end{merksatz}
